% Lampiran B: data lengkap. Tabel yang lebih lebar dari area teks
% dipindahkan ke lampiran (Pedoman Subbab 11.8 butir 4); contoh berikut
% memakai halaman landscape.
\lampiran[lamp:data]{Data Lengkap Hasil Kuesioner}

\panduan{Lampiran memuat materi pendukung yang terlalu panjang untuk
tubuh naskah, misalnya instrumen, surat izin mitra yang telah
dianonimkan, tabel hasil lengkap, log pencarian pustaka, dan kode sumber
utama (Subbab 11.2.4). Identitas responden diganti kode.}

Tabel~\ref{tab:data-lengkap} memuat contoh format data per responden.

\begin{landscape}
\begin{table}[htbp]
  \centering
  \caption{Contoh format data jawaban per responden (angka ilustratif)}
  \label{tab:data-lengkap}
  \begin{tabular}{@{}lccccccccccr@{}}
    \toprule
    Responden & B1 & B2 & B3 & B4 & B5 & B6 & B7 & B8 & B9 & B10 & $S_i$ \\
    \midrule
    R01 & 4 & 2 & 4 & 1 & 5 & 2 & 4 & 2 & 4 & 2 & 77,5 \\
    R02 & 3 & 3 & 4 & 2 & 4 & 2 & 3 & 3 & 4 & 3 & 62,5 \\
    R03 & 5 & 1 & 5 & 2 & 4 & 1 & 5 & 2 & 4 & 2 & 87,5 \\
    \bottomrule
  \end{tabular}
  \sumber{olahan penulis}
\end{table}
\end{landscape}
